% Moved verbatim from Paper A (new/sections/core.tex lines 300-331, commit f96523f) on 2026-10-29. Raw material for Paper B; rewrite, do not reuse.
\subsection{Loops}\label{sec:loops}
A loop of the host program runs when the circuit is built. A static slicer unrolls it, so every iteration
contributes its own operations, each tagged with the line and iteration that emit it. When the number of iterations
cannot be known statically, the loop body is repeated as \emph{may-operations}, operations that might be emitted,
until no repetition adds a location to any dependence chain. A loop of the circuit itself, Qiskit's
\texttt{while\_loop}, \texttt{for\_loop} or \texttt{switch}, runs on the device, and its condition reads measured
bits. A \texttt{for\_loop} over a known index set is unrolled, and any other circuit loop is repeated until stable
in the same way, with every operation in its body conditioned on the bits its condition reads, so the classical
control dependence on the measurement that decides the loop is kept.

\begin{proposition}[Loops]\label{prop:loops}
Let a loop body $B$ act on a set of locations $\Lambda$, let the loop end with probability one, and let $P_B$ repeat
$B$ $|\Lambda| + 1$ times, with every operation of $B$ conditioned on the bits that decide the loop. For a
criterion $W$, deleting from $B$ every operation none of whose copies is in the QDG slice of $P_B$ for $W$ preserves
$W$ at every input.
\end{proposition}

\begin{IEEEproof}[Proof sketch]
A loop that runs at most $K$ times is the circuit of $K$ copies of $B$, each conditioned on the bits that decide
the loop, a circuit with conditioned operations to which Theorem~\ref{thm:sound} applies. In its backward pass every
copy is processed alike, and the needed set only grows. Once a copy adds no location of $\Lambda$, every earlier copy
keeps the same operations, and the needed set can grow at most $|\Lambda|$ times, so for every $K$ the copies keep
only operations that one of the $|\Lambda| + 1$ copies of $P_B$ keeps. The bits that decide the loop are read by
every copy, so the measurements that write them are kept and the distribution over numbers of iterations is
preserved. A loop that ends with probability one is the limit of its bounded prefixes.
\end{IEEEproof}

The induction of Theorem~\ref{thm:sound} is not mechanized for loops, and the proposition is argued on paper. It is
checked exactly on \numCfPrograms{} programs of our own with circuit loops, branches and switches. For every one of
their \numCfQubits{} qubits, deleting every operation outside the qubit's slice from the source leaves the qubit's
state unchanged at \numCfInputs{} inputs, while the slices that read a circuit loop's body once are wrong on
\numCfOldWrong{} of the \numCfPrograms{} programs, an exploratory check of Section~\ref{sec:design}.
